\chapter{Hardware Security Analysis of Post-Quantum Cryptography}
\label{ch:hardware-security}

\abstract*{Deploying post-quantum cryptographic algorithms in hardware introduces attack surfaces that have no analogue in software implementations.  Every logic gate transition leaks information through power consumption and electromagnetic emanation; every register holds secrets that persist until actively destroyed.  This chapter develops hardware security analysis from first principles---the physics of side-channel leakage, the mathematics of masking countermeasures, and the engineering of fault-resilient architectures---then applies the full framework to ML-KEM, ML-DSA, and SLH-DSA.  We map the attack surfaces of each algorithm to specific countermeasure techniques with concrete cost estimates, analyse the cross-component interactions that make system-level security harder than the sum of its parts, and survey the certification landscape that governs whether a hardened implementation is accepted for deployment.}

\abstract{Deploying post-quantum cryptographic algorithms in hardware introduces attack surfaces that have no analogue in software implementations.  Every logic gate transition leaks information through power consumption and electromagnetic emanation; every register holds secrets that persist until actively destroyed.  This chapter develops hardware security analysis from first principles---the physics of side-channel leakage, the mathematics of masking countermeasures, and the engineering of fault-resilient architectures---then applies the full framework to ML-KEM, ML-DSA, and SLH-DSA.  We map the attack surfaces of each algorithm to specific countermeasure techniques with concrete cost estimates, analyse the cross-component interactions that make system-level security harder than the sum of its parts, and survey the certification landscape that governs whether a hardened implementation is accepted for deployment.}

% =============================================================
% Source material: generalised PQC hardware security analysis.
% Covers entropy-to-crypto pipeline, SCA/FI, masking, certification.
% =============================================================


\section{The Entropy-to-Cryptography Pipeline}
\label{sec:entropy-pipeline}

The previous chapter treated PQC algorithms as software: sequences of arithmetic instructions executed on a processor, where security means constant-time code and absence of secret-dependent branches.  Hardware changes the game entirely.  A dedicated PQC accelerator processes secrets not as abstract variables but as physical voltages on wires, currents through transistors, and charges stored in capacitors.  Every one of these physical quantities is, in principle, observable by an attacker with the right equipment.

The central object of this chapter is the \emph{entropy-to-cryptography pipeline}\index{entropy-to-cryptography pipeline}: the complete data path from the entropy source that generates randomness, through the deterministic random bit generator (DRBG) that conditions and expands it, to the PQC algorithms that consume it for key generation, encapsulation, and signing.  In a typical hardware PQC implementation, this pipeline has three stages:

\begin{enumerate}
\item \textbf{Entropy source.}\index{entropy source}  A physical random number generator (often a quantum random number generator, QRNG) that produces raw random bits.  These bits have high but imperfect min-entropy and must pass health tests before being trusted.
\item \textbf{DRBG.}\index{DRBG}  A cryptographic construction---typically based on SHAKE-256 or a similar extendable-output function---that absorbs entropy and produces the uniform random streams needed by PQC algorithms.  Its internal state (1600 bits for a Keccak-based DRBG) is the most valuable secret in the system: compromising it yields every key, nonce, and mask generated until the next reseed.
\item \textbf{PQC algorithms.}  ML-KEM\index{ML-KEM}, ML-DSA\index{ML-DSA}, and SLH-DSA\index{SLH-DSA}, each consuming random bits from the DRBG and processing long-term secrets (private keys) and ephemeral secrets (nonces, masks).
\end{enumerate}

The DRBG is the single chokepoint\index{DRBG!chokepoint} through which all entropy passes.  An attacker who recovers the DRBG state does not need to attack any individual PQC algorithm---all outputs become predictable.  Conversely, perfectly securing the PQC algorithms is worthless if the DRBG leaks its state through a power trace.

Figure~\ref{fig:ch13-pipeline} illustrates this data flow.  Each box processes at least one secret, and each arrow represents an interface that must be protected for both integrity (corruption leads to security failure) and confidentiality (observation leads to entropy leakage).

\begin{figure}[t]
\centering
\begin{tikzpicture}[
    block/.style={rectangle, draw, thick, minimum width=2.8cm, minimum height=1.2cm, align=center, font=\small},
    secret/.style={font=\footnotesize\itshape, text=red!70!black},
    >=Stealth,
    node distance=1.5cm and 2cm
]
    % Entropy source
    \node[block, fill=blue!10] (qrng) {Entropy Source\\(QRNG)};

    % DRBG
    \node[block, fill=orange!15, right=2.5cm of qrng] (drbg) {DRBG\\(XDRBG / SHAKE)};

    % PQC algorithms
    \node[block, fill=green!10, right=2.5cm of drbg, yshift=1.5cm] (mlkem) {ML-KEM\\(FIPS 203)};
    \node[block, fill=green!10, right=2.5cm of drbg] (mldsa) {ML-DSA\\(FIPS 204)};
    \node[block, fill=green!10, right=2.5cm of drbg, yshift=-1.5cm] (slhdsa) {SLH-DSA\\(FIPS 205)};

    % Arrows
    \draw[->, very thick] (qrng) -- node[above, font=\footnotesize] {raw bits} (drbg);
    \draw[->, thick] (drbg) -- (mlkem);
    \draw[->, thick] (drbg) -- (mldsa);
    \draw[->, thick] (drbg) -- (slhdsa);

    % Masking randomness bypass
    \draw[->, thick, dashed, red!60!black] (qrng) -- ++(0,-2) -| node[below, near start, font=\footnotesize, text=red!60!black] {masking randomness (raw)} (drbg.south);

    % Secret labels
    \node[secret, below=0.1cm of qrng] {quantum signal};
    \node[secret, below=0.1cm of drbg] {1600-bit state};
    \node[secret, above=0.1cm of mlkem] {$\hat{\mathbf{s}}$, seed $z$};
    \node[secret, below=0.1cm of mldsa] {$\mathbf{s}_1, \mathbf{s}_2, \mathbf{t}_0$};
    \node[secret, below=0.1cm of slhdsa] {\texttt{sk\_seed}};
\end{tikzpicture}
\caption{The entropy-to-cryptography pipeline.  Solid arrows carry conditioned randomness; the dashed arrow carries raw entropy directly to the DRBG's masking logic, breaking the circular dependency discussed in Sect.~\ref{sec:cross-component}.  Red labels indicate the primary secrets at each stage.}
\label{fig:ch13-pipeline}
\end{figure}

\begin{svgraybox}
\textbf{The weakest-link principle.}  In a hardware cryptographic module, the system security is determined by the least protected component in the entropy-to-cryptography pipeline.  Hardening ML-KEM to withstand $10^6$ attack traces is futile if the DRBG that generates ML-KEM's nonces can be broken with $10^3$ traces of unmasked Keccak.  Security analysis must be holistic: every component, every interface, every shared resource.
\end{svgraybox}

Table~\ref{tab:pipeline-components} summarises the attack surface of each pipeline component.  The DRBG stands out: its internal state is a long-lived secret processed repeatedly by Keccak, and without masking, CPA recovers the state in $10^2$--$10^4$ traces~\cite{BagheriEtAl2015}.  Multiple invocations between reseeds give the attacker correlated traces of the same secret, making the DRBG a more attractive target than any individual PQC algorithm.

\begin{table}[t]
\caption{Pipeline components, their secrets, and primary attack surfaces}
\label{tab:pipeline-components}
\begin{tabular}{p{2cm}p{3cm}p{3cm}p{4cm}}
\hline\noalign{\smallskip}
\textbf{Component} & \textbf{Primary Secret} & \textbf{Key Operations} & \textbf{Primary Attack Surface} \\
\noalign{\smallskip}\svhline\noalign{\smallskip}
Entropy source & Raw quantum signal & Detection, ADC, health tests & Physical manipulation, EM emanation \\
DRBG & Internal state (1600 bits) & Keccak-$f[1600]$ & CPA on Keccak, DFA on last rounds \\
ML-KEM & $\hat{\mathbf{s}}$, seed $z$ & NTT, compress & CPA on NTT, FI on comparison \\
ML-DSA & $\mathbf{s}_1, \mathbf{s}_2, \mathbf{t}_0$ & NTT, rejection sampling & Timing channel, CPA on Keccak \\
SLH-DSA & \texttt{sk\_seed} & Keccak (${\sim}35{,}000\times$) & CPA on WOTS+ chains \\
\noalign{\smallskip}\hline
\end{tabular}
\end{table}

This chapter analyses each stage of the pipeline, but the emphasis falls on the aspects unique to PQC.  Classical hardware security techniques (power analysis of AES, fault attacks on RSA-CRT) are well-documented~\cite{MangardOswaldPopp2007}; here we focus on the new attack surfaces that arise from lattice-based and hash-based constructions---the NTT, rejection sampling, polynomial compression, WOTS+ chains---and on the system-level interactions that make PQC hardware security qualitatively different from securing a single AES core.

The remainder of the chapter is organised as follows.  Section~\ref{sec:sca-first-principles} develops side-channel analysis from the physics of CMOS circuits through to the statistics of Correlation Power Analysis.  Section~\ref{sec:countermeasure-foundations} introduces the countermeasure toolkit: masking (Boolean and arithmetic), shuffling, fault injection defences, and zeroisation.  Sections~\ref{sec:mlkem-hardening}--\ref{sec:slhdsa-hardening} apply this toolkit to each of the three NIST-standardised PQC algorithms.  Section~\ref{sec:cross-component} examines the cross-component interactions that only manifest at the system level.  Section~\ref{sec:certification} surveys the certification frameworks that determine whether a hardened implementation is accepted for regulated deployment.


\section{Side-Channel Analysis from First Principles}
\label{sec:sca-first-principles}

A cryptographic algorithm, viewed as mathematics, processes secrets without any physical trace.  A cryptographic \emph{implementation}, viewed as a circuit, processes secrets by moving electrons through transistors---and every electron movement is, in principle, observable.  The gap between the mathematical abstraction and the physical reality is the side channel\index{side channel}.

\subsection{The Physics of Leakage}
\label{subsec:physics-leakage}

Consider a single CMOS inverter: a p-type transistor connected to the power supply ($V_{DD}$) and an n-type transistor connected to ground.  When the input transitions from 0 to 1, the n-transistor turns on and discharges the output capacitance to ground; when the input transitions from 1 to 0, the p-transistor turns on and charges the output capacitance from $V_{DD}$.  Each transition draws a transient current from the power supply.  A transition from 0 to 1 at the input produces a measurable current pulse; a 1 to 1 transition produces none.

Scale this to a register holding $n$ bits.  When the register is clocked to update from a value $v_{\text{old}}$ to a value $v_{\text{new}}$, the number of bit positions that change is the \emph{Hamming distance}\index{Hamming distance} $\text{HD}(v_{\text{old}}, v_{\text{new}}) = \text{HW}(v_{\text{old}} \oplus v_{\text{new}})$, where $\text{HW}$ denotes the Hamming weight\index{Hamming weight} (number of 1-bits).  Each changing bit charges or discharges a parasitic capacitance $C$, drawing an energy of approximately $\frac{1}{2}C V_{DD}^2$ per transition.  The total power consumed in that clock cycle is therefore approximately proportional to $\text{HD}(v_{\text{old}}, v_{\text{new}})$, plus a data-independent component (clock tree, leakage current, other logic).

This is the \emph{Hamming distance power model}\index{Hamming distance model}~\cite{MangardOswaldPopp2007}:
\begin{equation}
\label{eq:hd-model}
P(t) \;\approx\; \alpha \cdot \text{HD}(v_{\text{old}}, v_{\text{new}}) + \beta + N(t),
\end{equation}
where $\alpha$ is a gain factor depending on the capacitance and supply voltage, $\beta$ captures the data-independent power, and $N(t)$ is measurement noise.  The same physics produces electromagnetic emanation: each current pulse generates a transient magnetic field that can be measured with a near-field EM probe, sometimes with better spatial resolution than power measurements.

\begin{svgraybox}
\textbf{The fundamental problem.}  If $v_{\text{old}}$ is known (or constant, such as an initialisation value) and $v_{\text{new}}$ depends on a secret $k$, then measuring $P(t)$ leaks information about $k$.  No amount of algorithmic cleverness can eliminate this leakage---it is a consequence of the laws of physics applied to digital circuits.  The only options are to make the leakage statistically undetectable (masking) or to reduce the number of observations the attacker can collect (limiting trace availability).
\end{svgraybox}


\subsection{What a Power Trace Looks Like}
\label{subsec:power-trace}

An attacker places a small resistor (1--10~$\Omega$) in series with the target device's power supply and measures the voltage drop across it with a digital oscilloscope sampling at 1--5~GS/s.  The result is a \emph{power trace}\index{power trace}: a time series of voltage samples, one per nanosecond or better, spanning the duration of the cryptographic operation.  For a Keccak permutation running at 100~MHz, a single invocation takes 24 rounds $\times$ 1 cycle/round $= 240$~ns, producing roughly 1000 samples at 4~GS/s.

Each clock cycle appears as a spike in the trace, with amplitude modulated by the data being processed.  The trace of an ML-KEM decapsulation at 100~MHz might span 50~$\mu$s and contain $200{,}000$ samples.  Hidden in those samples is enough information to recover the secret key---if the attacker knows where to look and how to extract the signal from the noise.


\subsection{Simple Power Analysis}
\label{subsec:spa}

The most direct attack is \emph{Simple Power Analysis} (SPA)\index{Simple Power Analysis}\index{SPA}: visually inspecting a single (or very few) power trace(s) to identify operations or data values.  In RSA, the square-and-multiply pattern reveals the exponent bits directly from one trace.  In PQC, SPA targets are less dramatic but still real: the number of iterations in ML-DSA's rejection sampling loop is visible as repeated blocks in the trace, and the accept/reject decision may produce a distinguishable power signature.

SPA requires no statistics---just a good oscilloscope and knowledge of the algorithm's structure.  It is the reason that every conditional branch on a secret must be implemented as a constant-time, data-independent operation: not merely because of timing differences, but because the power profile of taken versus not-taken branches differs measurably.

In PQC, SPA is less immediately devastating than in RSA (where the entire private key can be read from one trace), but it still matters.  The variable-time rejection loop in ML-DSA signing is an SPA-class vulnerability: the number of loop iterations is visible in the trace as a repeated pattern of NTT, multiplication, and norm-check blocks.  Even if the per-iteration power profile does not directly reveal secret coefficients, the iteration count leaks information about the secret key over many signatures (Sect.~\ref{subsec:rejection-sampling}).


\subsection{Differential and Correlation Power Analysis}
\label{subsec:dpa-cpa}

When a single trace does not reveal the secret directly, the attacker turns to statistics.  \emph{Differential Power Analysis} (DPA)\index{Differential Power Analysis}\index{DPA}, introduced by Kocher, Jaffe, and Jun~\cite{KocherJaffeJun1999}, and its refined successor \emph{Correlation Power Analysis} (CPA)\index{Correlation Power Analysis}\index{CPA}~\cite{BrierClavier2004}, exploit the statistical relationship between the secret and the measured power across many executions of the same operation with different inputs.

The CPA procedure, explained step by step, proceeds as follows.

\begin{definition}[Correlation Power Analysis]
\label{def:cpa}
Let $k^*$ be the secret (e.g., one coefficient of the ML-KEM secret polynomial $\hat{s}$), and let $x_i$ be the known input for the $i$-th trace.  The attacker:
\begin{enumerate}
\item Collects $N$ power traces $\{P_i(t)\}_{i=1}^N$, each corresponding to the processing of input $x_i$ under the same secret $k^*$.
\item For each key hypothesis $k \in \mathcal{K}$ (e.g., $k \in \{0, 1, \ldots, q-1\}$ for a single coefficient mod~$q$):
  \begin{enumerate}
  \item[(a)] Computes the hypothetical intermediate value $v_i(k) = f(x_i, k)$ for each trace, where $f$ is a known function (e.g., a modular multiplication in the NTT butterfly).
  \item[(b)] Predicts the power consumption $h_i(k) = \text{HW}(v_i(k))$ or $\text{HD}(v_i^{\text{prev}}, v_i(k))$ using the Hamming weight or Hamming distance model.
  \item[(c)] Computes the Pearson correlation coefficient between the vector $(h_1(k), \ldots, h_N(k))$ and the measured power at each time sample:
  \begin{equation}
  \label{eq:cpa-correlation}
  \rho_k(t) = \frac{\sum_{i=1}^N (h_i(k) - \bar{h}_k)(P_i(t) - \bar{P}(t))}{\sqrt{\sum_{i=1}^N (h_i(k) - \bar{h}_k)^2 \cdot \sum_{i=1}^N (P_i(t) - \bar{P}(t))^2}}.
  \end{equation}
  \end{enumerate}
\item Selects the hypothesis $k$ that achieves the highest absolute correlation: $\hat{k} = \arg\max_k \max_t |\rho_k(t)|$.
\end{enumerate}
If $N$ is large enough, $\hat{k} = k^*$ with overwhelming probability.
\end{definition}

The intuition behind CPA is worth pausing over, because it is the foundation of all statistical side-channel attacks.  For the correct key hypothesis $k = k^*$, the predicted power $h_i(k^*)$ is \emph{genuinely correlated} with the measured power at the relevant time sample: both reflect the same underlying data value.  For any incorrect hypothesis $k \neq k^*$, the predicted power is a function of a \emph{wrong} intermediate value and therefore uncorrelated with the measurement.  As $N$ grows, the correlation for the correct hypothesis converges to a non-zero value (determined by the model accuracy and the SNR), while the correlations for incorrect hypotheses converge to zero by the law of large numbers.  The attack succeeds when the gap between the correct peak and the tallest incorrect peak becomes statistically significant.

The critical question is: how large must $N$ be?  For an unprotected hardware implementation of AES, the answer is 50--500 traces.  For an unprotected ML-KEM NTT, published attacks demonstrate full key recovery in $10^3$--$10^4$ traces~\cite{RaviRoyBhasin2020, NgoRaviSinhaRoy2021}.  The ML-KEM trace count is higher than AES for two reasons: the modulus $q = 3329$ means each coefficient has $3329$ possible values (versus $256$ for an AES byte), spreading the correlation signal across more hypotheses; and the NTT butterfly involves multiple operations per coefficient, blurring the temporal localisation.  The exact number depends on the signal-to-noise ratio (SNR): a dedicated ASIC with low switching noise requires fewer traces than an FPGA sharing its power supply with other logic.

A useful rule of thumb relates the SNR to the trace count.  For a first-order attack (no masking), the minimum number of traces scales as:
\begin{equation}
\label{eq:trace-count-snr}
N \;\approx\; \frac{\alpha}{\text{SNR}^2},
\end{equation}
where $\alpha$ is a constant depending on the model accuracy and the number of hypotheses (typically $\alpha \approx 3$--$10$).  For a first-order masked implementation, the attacker must perform a \emph{second-order} analysis---correlating products or combinations of leakage from two different time samples---and the trace count scales as $N \approx \alpha / \text{SNR}^4$.  Each additional masking order squares the exponent, explaining the exponential security gain with masking order.

\begin{example}[CPA on an NTT Butterfly]
\label{ex:cpa-ntt}
Consider a single NTT butterfly computing $a' = a + \omega \cdot b \bmod q$ where $\omega$ is a known twiddle factor and $b$ is a secret coefficient of $\hat{s}$.  The attacker knows $a$ (it is derived from the public ciphertext) and $\omega$ (it is a fixed constant).  For each hypothesis $b = k$, the attacker predicts $v(k) = a + \omega \cdot k \bmod q$ and models the power as $\text{HW}(v(k))$.  After collecting traces from 5000 decapsulations (with different ciphertexts but the same secret key), the correlation peak for the correct $k = b$ rises above the noise, while all incorrect hypotheses remain at the noise floor.  Repeating for each of the 256 coefficients recovers the full secret polynomial.

To make this concrete with numbers: for $q = 3329$, $\omega = 17$, and $a = 1000$, the intermediate value for $k = 42$ is $v(42) = 1000 + 17 \times 42 \bmod 3329 = 1714$, with $\text{HW}(1714) = \text{HW}(\texttt{0x6B2}) = 7$.  For $k = 43$, $v(43) = 1731$, with $\text{HW}(1731) = 8$.  Different key hypotheses produce different Hamming weights, and the correlation distinguishes the correct one---provided enough traces are averaged to suppress the noise.
\end{example}


\subsection{Electromagnetic Emanation}
\label{subsec:em-emanation}

Power analysis measures the \emph{aggregate} current drawn by the entire chip.  Electromagnetic (EM) analysis\index{electromagnetic analysis} measures the \emph{local} magnetic field above a specific region of the die, using a near-field probe with a loop diameter of 100--500~$\mu$m.  The physics is identical---current flow creates magnetic fields---but EM probes offer two advantages: spatial selectivity (targeting the specific register that processes the secret, ignoring noise from unrelated logic) and the ability to bypass on-chip power filtering that attenuates the power-supply signal.

For FPGA implementations, where the placement of logic is partially under the designer's control, EM probes can isolate the Keccak core, the NTT multiplier, or the comparison logic.  The cost of a professional EM measurement setup ranges from \$50K to \$500K, compared to \$5K--\$20K for a power analysis setup.  However, the information gained per trace is often higher, reducing the total number of traces needed.

The practical implication for PQC hardware designers is that power-rail filtering and decoupling capacitors---standard practice for signal integrity---provide some attenuation of the power-supply signal but offer no protection against EM probes.  A device that ``passes'' a power-analysis test may still fail an EM test.  Certification laboratories (Sect.~\ref{sec:certification}) typically test both channels.


\subsection{Microarchitectural Side Channels}
\label{subsec:cache-timing}

When PQC algorithms run on general-purpose processors rather than dedicated hardware, a different class of side channels emerges.  \emph{Cache-timing attacks}\index{cache-timing attacks} exploit the fact that memory access time depends on whether the data is in the cache (fast) or main memory (slow).  If the secret determines which memory addresses are accessed---as in a lookup table indexed by a secret value---the attacker can infer the secret by measuring access times or by priming and probing cache lines.

These attacks are devastating against table-based implementations of Keccak or polynomial sampling, but they are primarily a software concern.  In dedicated hardware, there is no cache.  We mention them here for completeness, because many ``hardware PQC'' deployments actually use a soft-core processor running PQC software on an FPGA, inheriting all the software side channels plus the hardware ones.

The distinction between ``hardware PQC'' and ``software PQC on hardware'' is crucial for security analysis.  A soft-core RISC-V processor on an FPGA running the reference C implementation of ML-KEM is vulnerable to cache-timing attacks, branch-prediction attacks, and power analysis of the processor's ALU---simultaneously.  A dedicated RTL implementation of ML-KEM on the same FPGA has none of the microarchitectural side channels but is fully exposed to power and EM analysis of its custom datapath.  The threat models are fundamentally different, and countermeasures that work for one may be irrelevant for the other.  This chapter focuses on dedicated hardware implementations; Chapter~12 addressed the software case.


\subsection{Threat Model Summary}
\label{subsec:threat-model}

\begin{table}[t]
\caption{Side-channel attack classes, equipment cost, and typical trace requirements for unprotected PQC implementations}
\label{tab:sca-classes}
\begin{tabular}{p{3cm}p{2.5cm}p{2.5cm}p{4cm}}
\hline\noalign{\smallskip}
\textbf{Attack Class} & \textbf{Equipment Cost} & \textbf{Traces Needed} & \textbf{Primary Target} \\
\noalign{\smallskip}\svhline\noalign{\smallskip}
SPA               & \$5K--\$20K   & 1--10        & Control flow, iteration count \\
DPA/CPA (power)   & \$5K--\$20K   & $10^3$--$10^4$ & NTT coefficients, Keccak state \\
CPA (EM)          & \$50K--\$500K & $10^2$--$10^3$ & Targeted registers \\
Cache timing      & \$0 (software) & $10^2$--$10^6$ & Table lookups, branches \\
Fault injection   & \$10K--\$100K & 1--10        & Comparison logic, loop counters \\
\noalign{\smallskip}\hline
\end{tabular}
\end{table}

Table~\ref{tab:sca-classes} summarises the landscape.  The attacker's capability is not exotic: a \$10K oscilloscope and physical access to the device suffice for power-based CPA.  Certification laboratories (Sect.~\ref{sec:certification}) routinely perform these attacks as part of the evaluation process.  The remainder of this chapter is devoted to the countermeasures that make these attacks fail---or at least raise their cost beyond practical feasibility.


\section{Countermeasure Foundations}
\label{sec:countermeasure-foundations}

The previous section established that side-channel leakage is a physical inevitability.  This section develops the mathematical and engineering tools for rendering that leakage useless to an attacker.  The central technique is \emph{masking}\index{masking}: splitting every secret into random shares such that no single share reveals anything about the original.  But masking alone is not enough---we also need defences against fault injection, techniques for secure memory management, and a clear understanding of where each countermeasure applies within the PQC pipeline.


\subsection{Boolean Masking}
\label{subsec:boolean-masking}

The simplest form of masking splits a secret $x$ into two shares\index{masking!Boolean} $x_0$ and $x_1$ such that $x = x_0 \oplus x_1$, where $x_0$ is chosen uniformly at random.  The share $x_1 = x \oplus x_0$ is then also uniformly random (since XOR with a uniform value produces a uniform result).  An attacker who observes only $x_0$ or only $x_1$---through a side channel that captures one register at a time---learns nothing about $x$.  Recovering $x$ requires combining information from \emph{both} shares simultaneously, which requires a qualitatively different (and much harder) attack.

\begin{example}[Boolean Masking of a Single Byte]
\label{ex:boolean-masking}
Let $x = \texttt{0xA7}$ (binary \texttt{10100111}).  We choose a uniform random mask $x_0 = \texttt{0x3C}$ (binary \texttt{00111100}).  The second share is $x_1 = x \oplus x_0 = \texttt{0x9B}$ (binary \texttt{10011011}).  The Hamming weight of $x$ is 5, but the Hamming weights of $x_0$ and $x_1$ are 4 and 5 respectively---both unrelated to the secret.  An attacker measuring the power consumed when processing $x_0$ or $x_1$ sees a value uniformly distributed over all 256 possible bytes, regardless of $x$.  To recover $x$, the attacker would need to observe $x_0$ and $x_1$ in the same clock cycle (or in two cycles and combine them statistically), which requires a second-order attack.
\end{example}

More generally, $d$-th order Boolean masking splits $x$ into $d+1$ shares:
\begin{equation}
\label{eq:boolean-masking}
x = x_0 \oplus x_1 \oplus \cdots \oplus x_d,
\end{equation}
where $x_0, \ldots, x_{d-1}$ are independent uniform random values and $x_d = x \oplus x_0 \oplus \cdots \oplus x_{d-1}$.  A $d$-th order attack---one that combines leakage from $d+1$ time samples simultaneously---is needed to extract information about $x$.  The number of traces required for a $d$-th order attack scales roughly as $\text{SNR}^{-2d}$, where $\text{SNR}$ is the signal-to-noise ratio~\cite{MangardOswaldPopp2007}.  First-order masking ($d=1$) is the practical minimum for certification; it squares the trace requirement compared to an unmasked implementation.

The beauty of Boolean masking is that \emph{linear operations are free}.  If we want to compute $z = x \oplus y$ where both $x$ and $y$ are masked, we simply XOR corresponding shares: $z_0 = x_0 \oplus y_0$, $z_1 = x_1 \oplus y_1$.  No fresh randomness is needed, and each share computation touches only one share of each input.

The difficulty arises with \emph{non-linear operations}.  Computing $z = x \wedge y$ (bitwise AND) from shares requires evaluating cross-terms: $(x_0 \oplus x_1) \wedge (y_0 \oplus y_1) = (x_0 \wedge y_0) \oplus (x_0 \wedge y_1) \oplus (x_1 \wedge y_0) \oplus (x_1 \wedge y_1)$.  The cross-terms $x_0 \wedge y_1$ and $x_1 \wedge y_0$ each depend on shares of different secrets, and computing them in the same combinational circuit creates a transient moment where both shares influence the same wire---exactly the leakage that masking was supposed to prevent.

\begin{definition}[ISW Scheme]
\label{def:isw}
\index{ISW scheme}
The Ishai--Sahai--Wagner (ISW) scheme~\cite{IshaiSahaiWagner2003} computes a masked AND by introducing \emph{fresh randomness} to break the correlation between cross-terms.  For first-order masking ($d=1$), the masked AND of $(x_0, x_1)$ and $(y_0, y_1)$ proceeds as:
\begin{enumerate}
\item Draw a fresh random bit $r$.
\item Compute $z_0 = (x_0 \wedge y_0) \oplus r \oplus (x_0 \wedge y_1) \oplus (x_1 \wedge y_0)$.
\item Compute $z_1 = (x_1 \wedge y_1) \oplus r$.
\end{enumerate}
Each intermediate computation depends on at most one share of each input, plus fresh randomness.  The cost: one fresh random bit per AND gate.
\end{definition}

In practice, the ISW scheme alone is insufficient for hardware implementations.  The reason is \emph{glitches}\index{glitches}: in combinational logic, signals do not arrive at a gate simultaneously.  If $x_0 \wedge y_1$ is computed slightly before $r$ is XORed in, there is a brief moment where the wire carries a value that depends on both $x_0$ and $y_1$---and that glitch is observable in the power trace.


\subsection{Domain-Oriented Masking}
\label{subsec:dom}

\emph{Domain-Oriented Masking} (DOM)\index{Domain-Oriented Masking}\index{DOM}~\cite{GrossHankeDeMulder2016} solves the glitch problem by inserting a \emph{pipeline register} between the computation of cross-domain terms and their recombination.  The register ensures that each share is stable for a full clock cycle before it interacts with shares from another domain, preventing transient glitch-based leakage.

Why do glitches matter so much?  In a hardware implementation, combinational logic does not compute instantaneously.  When inputs to a gate arrive at different times (due to different propagation delays through preceding logic), the gate output may transition through intermediate values before settling.  If a masked AND gate receives $x_0$ before $r$ (the fresh random bit), there is a transient nanosecond-scale window where the output carries $x_0 \wedge y_1$---a cross-domain term that depends on shares from both the $x$ and $y$ secrets.  This transient draws current from the power supply, and the current is observable in the power trace.  The ISW scheme is provably secure in the \emph{probing model} (where the attacker reads individual wires), but not in the \emph{glitch-extended probing model} that captures this physical reality.  DOM closes this gap.

For first-order DOM (DOM-1), the overhead per AND gate is one fresh random bit plus one pipeline cycle of latency.  For a complex function like the Keccak permutation, the overhead concentrates in the single non-linear step: the $\chi$ function.  All other steps of Keccak ($\theta$, $\rho$, $\pi$, $\iota$) are linear (XOR-based) and apply independently to each share at zero cost.  The $\chi$ step computes $a[i] = a[i] \oplus ((\neg a[i{+}1]) \wedge a[i{+}2])$---one AND per bit of the 1600-bit state.  DOM-1 on $\chi$ requires $1600$ fresh random bits per round $\times$ $24$ rounds $= 38{,}400$ fresh random bits per Keccak-$f[1600]$ invocation.

At a Keccak throughput of 100 invocations per second, the masking randomness demand is approximately 3.84~Mbit/s---a non-trivial but achievable requirement for a hardware random number generator.


\subsection{Arithmetic Masking}
\label{subsec:arithmetic-masking}

Boolean masking works naturally with XOR-based operations (Keccak, bitwise comparisons).  But PQC algorithms spend much of their time in modular arithmetic: the NTT operates in $\Zq$, polynomial multiplication is coefficient-wise multiplication mod~$q$, and key generation involves additions mod~$q$.  For these operations, \emph{arithmetic masking}\index{masking!arithmetic} is the natural choice:
\begin{equation}
\label{eq:arithmetic-masking}
x = x_0 + x_1 \bmod q.
\end{equation}

Under arithmetic masking, all linear modular operations---addition, subtraction, multiplication by a \emph{public} constant---apply independently to each share at no extra cost.  If $c$ is a public constant and $x = x_0 + x_1 \bmod q$, then $c \cdot x = (c \cdot x_0) + (c \cdot x_1) \bmod q$.  No fresh randomness, no cross-terms.

This observation has a profound consequence for PQC.

\begin{svgraybox}
\textbf{The NTT is free under arithmetic masking.}\index{NTT!masking cost}  The NTT butterfly computes $a' = a + \omega b \bmod q$ and $b' = a - \omega b \bmod q$, where $\omega$ is a public twiddle factor.  Since $\omega$ is public, the multiplication $\omega \cdot b$ is a public-constant--times--secret operation: linear under arithmetic masking.  The additions and subtractions are also linear.  Therefore, the \emph{entire} NTT and pointwise polynomial multiplication (where the public matrix $\mathbf{A}$ multiplies the secret vector $\hat{\mathbf{s}}$) can be masked by simply running each share through the NTT independently.  Cost: $2\times$ the multiplier area, zero fresh randomness.
\end{svgraybox}

Where, then, does the masking cost concentrate?  At the \emph{boundaries} between arithmetic and Boolean domains.


\subsection{The A2B/B2A Conversion Problem}
\label{subsec:a2b-b2a}

PQC algorithms inevitably switch between arithmetic operations (NTT, modular multiply) and Boolean operations (Keccak for hashing, bitwise comparison for decapsulation checks, compression via rounding).  At each switch, the masked representation must be converted\index{masking!A2B conversion}\index{masking!B2A conversion}:

\begin{itemize}
\item \textbf{A2B (Arithmetic to Boolean):} Given shares $x_0, x_1$ with $x = x_0 + x_1 \bmod q$, produce shares $y_0, y_1$ with $x = y_0 \oplus y_1$.  This requires a \emph{masked carry chain}: the addition involves carries that propagate bit by bit, and each carry computation is a non-linear (AND) operation that must be masked.  Cost: $\bigO(n)$ masked ANDs for $n$-bit values.
\item \textbf{B2A (Boolean to Arithmetic):} The reverse conversion, typically implemented via Goubin's algorithm~\cite{Goubin2001} or optimised variants.
\end{itemize}

The concrete cost depends on the bit-width of $q$.  For ML-KEM with $q = 3329$ (12 bits): approximately 12 cycles and 12 bits of fresh randomness per conversion.  For ML-DSA with $q = 8380417$ (23 bits): approximately 23 cycles and 23 bits of fresh randomness.  The cost scaling with bit-width means ML-DSA masking is roughly twice as expensive per coefficient as ML-KEM masking---a ratio that pervades the entire hardening analysis.

The A2B/B2A conversion is not merely an overhead---it is a \emph{verified security boundary}\index{A2B conversion!security boundary}.  Each conversion must be proven to maintain the masking invariant: after the conversion, each output share must be statistically independent of the secret.  Errors in the conversion logic (incorrect handling of carry propagation, off-by-one in the modular reduction) are among the most subtle and dangerous implementation bugs in masked PQC hardware, because they silently break the security guarantee without producing any functional error.


\subsection{Shuffling}
\label{subsec:shuffling}

\emph{Shuffling}\index{shuffling} processes polynomial coefficients in a random permutation that changes with every execution.  If the attacker expects coefficient $i$ to appear at time step $i$, but the implementation processes coefficient $\pi(i)$ instead (for a random permutation $\pi$), the attacker's power model predictions are misaligned with the actual computation.  The number of traces required increases by a factor proportional to the number of possible positions---up to $256\times$ for a polynomial of 256 coefficients.

Shuffling is implemented via the Fisher--Yates algorithm in hardware (approximately 256 clock cycles of setup, consuming roughly 2~Kbits of randomness) or via a cheaper partial permutation (random start plus coprime stride, which offers fewer permutations but is simpler to implement).

A critical point: shuffling is a \emph{statistical} countermeasure, not a cryptographic one.  It increases the number of traces needed but does not provide information-theoretic protection.  It must always \emph{complement} masking, never replace it.

\begin{example}[Shuffling the NTT]
\label{ex:shuffling-ntt}
Consider an NTT on a polynomial of 256 coefficients.  Without shuffling, the attacker knows that coefficient $i$ is processed at butterfly stage $j$ at a specific clock cycle.  With shuffling, the processing order is a random permutation.  The attacker's CPA hypothesis for coefficient $i$ must be tested against all 256 possible time positions, and the correlation signal is diluted by a factor of $\sqrt{256} = 16$.  Combined with first-order masking, the effective trace requirement increases from roughly $10^6$ (masking alone) to roughly $10^6 \times 256 \approx 2.6 \times 10^8$ (masking plus shuffling).  This is well beyond practical feasibility for most attack scenarios.
\end{example}


\subsection{Fault Injection Countermeasures}
\label{subsec:fi-countermeasures}

Side-channel analysis is a \emph{passive} attack: the attacker observes but does not interfere.  \emph{Fault injection}\index{fault injection} (FI) is an \emph{active} attack: the attacker deliberately disrupts the computation---via voltage glitching, clock manipulation, laser pulses, or electromagnetic pulses---to induce errors that reveal secrets or bypass security checks.

Three classes of countermeasures address fault injection:

\textbf{Temporal redundancy.}\index{temporal redundancy}  Compute the operation twice and compare results before outputting.  If a transient fault corrupts one computation, the comparison fails and the device aborts.  Cost: $2\times$ time or, with parallel redundancy, $2\times$ area.  Temporal redundancy is effective against \emph{transient} faults (voltage glitches, single laser pulses) but not against \emph{permanent} faults (circuit modifications by a well-resourced attacker).  For the highest threat levels, two independent implementations on different die regions provide stronger guarantees.

\textbf{Information redundancy.}\index{information redundancy}  Maintain parity, error-correcting codes, or CRCs on intermediate values, verified at each pipeline stage.  A fault that modifies a secret value also corrupts the check bits with high probability, triggering detection.  The error-detection code must be chosen so that the fault model of interest (single-bit flips for voltage glitches, multi-bit burst errors for laser pulses) falls within the code's detection capability.

\textbf{FSM protection.}\index{FSM protection}  The control logic (finite state machine) that sequences cryptographic operations is itself a target.  One-hot encoding with illegal-state detection, or Hamming-coded states with a trap state that triggers immediate zeroisation, prevents an attacker from skipping critical steps (such as the re-encryption check in ML-KEM decapsulation or the norm check in ML-DSA signing).  The trap state is critical: when an illegal state is detected, the module must not simply reset (which might leave secrets in memory) but must actively zeroise all secret-holding locations before any further operation.

For PQC specifically, high-value FI targets include the comparison logic in ML-KEM's implicit rejection, the rejection loop counter in ML-DSA, and the address register in SLH-DSA's Merkle tree traversal.  We will examine each in the per-algorithm sections that follow.


\subsection{Zeroisation}
\label{subsec:zeroisation}

After a cryptographic operation completes, every register and memory location that held secret data must be actively overwritten with a known pattern (typically all zeros, then all ones, then all zeros again)\index{zeroisation}.  Simply deasserting a signal or relying on power-down to clear memory is insufficient: FPGA block RAMs (BRAMs) do not self-clear on reset, and register contents may persist through soft resets.

The zeroisation path must be \emph{independent} of the normal operational path---otherwise, a single fault injection that disrupts the normal completion sequence also bypasses cleanup.  FIPS 140-3 mandates zeroisation under SP 800-140D~\cite{FIPS140-3}; Common Criteria covers it under FCS\_CKM.4~\cite{CommonCriteria2022}.  Verification of zeroisation completeness---confirming that \emph{every} secret-holding location is covered---is one of the most tedious but most important steps in hardware security certification.

For PQC implementations, the zeroisation scope is larger than for classical algorithms.  An AES core has a 128-bit key register and a 128-bit state register---two locations to clear.  An ML-KEM decapsulation core has the secret polynomial $\hat{\mathbf{s}}$ (multiple kilobits across NTT butterfly buffers), the seed $z$, intermediate Keccak states, masked shares, and the comparison result---dozens of locations distributed across the datapath.  A systematic zeroisation analysis begins by enumerating every register and BRAM location that touches secret data at any point during the computation, then verifying that each location is overwritten on: (a) normal completion, (b) fault/tamper detection, and (c) power-down or reset.

\begin{svgraybox}
\textbf{FPGA BRAMs do not self-clear.}\index{BRAM!zeroisation}  Unlike flip-flops, which typically power up in a defined state, FPGA block RAMs may retain their contents across configuration reloads and soft resets.  A ``cold boot'' attack on an FPGA that stores ML-KEM secret keys in BRAM is realistic: removing power briefly and re-reading the BRAM may recover partial key material.  Active zeroisation on shutdown---triggered by a voltage-drop detector with enough energy storage to complete the overwrite sequence---is the countermeasure.
\end{svgraybox}

With the countermeasure toolkit in hand, we now apply it to each of the three NIST-standardised PQC algorithms, identifying the specific attack surfaces and the cost of hardening each one.


\section{ML-KEM: Attack Surfaces and Hardening}
\label{sec:mlkem-hardening}

ML-KEM\index{ML-KEM!hardware security} (FIPS 203~\cite{FIPS203}) operates over $\Zq$ with $q = 3329$ (12 bits).  Among the three PQC standards, it is the most amenable to hardware masking, because most of its arithmetic is linear under arithmetic masking and the modulus is small.  The primary concern is the \emph{decapsulation} operation, which processes the long-term secret key $\hat{\mathbf{s}}$ once per key exchange---potentially millions of times over the key's lifetime.

\subsection{Decapsulation: A Guided Tour of the Attack Surface}
\label{subsec:mlkem-decaps}

ML-KEM decapsulation proceeds through six stages, each with distinct security implications.  We walk through them in order, marking the secret data processed at each step and the applicable attack class.

\textbf{Step 1: NTT and pointwise multiplication.}  The ciphertext vector $\hat{\mathbf{u}}$ (public) is multiplied by the secret key $\hat{\mathbf{s}}$ in the NTT domain.  As established in Sect.~\ref{subsec:arithmetic-masking}, this is fully linear under arithmetic masking: each share of $\hat{\mathbf{s}}$ is multiplied by the public $\hat{\mathbf{u}}$ independently.  Without masking, CPA recovers the secret coefficients in $10^3$--$10^4$ traces~\cite{RaviRoyBhasin2020}.  With first-order arithmetic masking, the attack requires second-order statistical analysis, raising the trace requirement to roughly $10^6$--$10^8$ depending on SNR.

\textbf{Step 2: Inverse NTT and decompression.}  The inverse NTT is also linear (public twiddle factors operating on masked shares).  Decompression involves scaling by powers of 2, which are public constants---still linear.

\textbf{Step 3: CBD sampling.}\index{CBD sampling}  When re-encapsulating to verify the ciphertext (the Fujisaki--Okamoto transform), the secret polynomial is regenerated from a seed using the Centred Binomial Distribution.  This involves: (a) Keccak hashing of the secret seed, requiring DOM-1 Keccak; (b) a popcount operation on secret data, which is non-linear and requires Boolean masking.  The B2A conversion (from the Boolean domain of Keccak output to the arithmetic domain of polynomial operations) costs approximately 12 cycles per coefficient for $q = 3329$.

\textbf{Step 4: Compression.}\index{compression (ML-KEM)}  The polynomial is compressed by rounding: a non-linear operation that requires A2B conversion (from arithmetic masking to Boolean representation for the rounding comparison).  This is a localised but critical unmasking point.

\textbf{Step 5: Comparison (implicit rejection).}  The re-encrypted ciphertext $c'$ is compared bitwise against the received ciphertext $c$.  This comparison must be \emph{constant-time and constant-power}: the implementation must not reveal whether $c' = c$ or $c' \neq c$.  A distinguishable comparison gives the attacker a chosen-ciphertext oracle, breaking CCA security.  The standard approach is a ``balanced MUX'' that always computes both the real and rejection shared secrets and selects the output using a secret bit derived from the comparison result.

\textbf{Step 6: Shared secret derivation.}  Keccak processes the secret seed $z$ (used for implicit rejection) and the derived shared secret.  The seed $z$ must be masked throughout.

\begin{table}[t]
\caption{ML-KEM decapsulation: attack surfaces and countermeasures}
\label{tab:mlkem-surfaces}
\begin{tabular}{p{3.5cm}p{1.5cm}p{3cm}p{4cm}}
\hline\noalign{\smallskip}
\textbf{Operation} & \textbf{Risk} & \textbf{Attack Class} & \textbf{Countermeasure} \\
\noalign{\smallskip}\svhline\noalign{\smallskip}
NTT / pointwise multiply & Critical & CPA ($10^3$--$10^4$ traces) & Arithmetic masking (free) \\
SHAKE for secret sampling & Critical & CPA on Keccak state & DOM-1 Keccak \\
CBD popcount & High & SCA on non-linear op & Masked popcount, B2A \\
Compression rounding & Medium & SCA at unmasking point & A2B, late unmasking \\
Decaps comparison & High & Timing/power oracle & Constant-power balanced MUX \\
FI on comparison & High & CCA oracle & Temporal redundancy \\
Secret key in BRAM & Medium & Probing, cold boot & Masked storage, ECC \\
\noalign{\smallskip}\hline
\end{tabular}
\end{table}


\subsection{Masking Cost Summary}
\label{subsec:mlkem-masking-cost}

The cost structure of masking ML-KEM is lopsided in a fortunate way.  The most computationally intensive operations---the NTT and pointwise multiplication---are free under arithmetic masking (aside from the $2\times$ multiplier area for two shares).  The real cost concentrates at the domain boundaries:

\begin{itemize}
\item \textbf{CBD sampling:} B2A conversion, $\sim$$256 \times 12 = 3{,}072$ cycles per polynomial.
\item \textbf{Compression:} A2B conversion, $\sim$$256 \times 12 = 3{,}072$ cycles per polynomial.
\item \textbf{Keccak:} DOM-1 adds one pipeline cycle per round (24 cycles per invocation) and demands 38,400 bits of fresh randomness per invocation.
\item \textbf{Comparison + MUX:} A2B for bitwise compare, plus balanced MUX logic.
\end{itemize}

For a hardware implementation running at 100~MHz, the masking overhead adds roughly 20--30\% to the decapsulation latency---dominated by the Keccak DOM-1 cost and the A2B/B2A conversions, not by the NTT.

\subsection{The Implicit Rejection Mechanism}
\label{subsec:implicit-rejection}

ML-KEM's CCA security relies on the Fujisaki--Okamoto transform~\cite{HofheinzHovemannKiltz2017}: after decrypting, the module re-encrypts the plaintext and compares the result with the received ciphertext.  If they match, the real shared secret is returned; if not, a pseudorandom ``rejection'' shared secret (derived from the secret seed $z$ and the ciphertext) is returned instead.  The attacker must not be able to distinguish which case occurred.

In hardware, this means: (1) both the real and rejection shared secrets must \emph{always} be computed, regardless of the comparison result; (2) the multiplexer that selects between them must draw the same power in both cases (a \emph{balanced MUX}); (3) the comparison must evaluate all bytes without early termination; and (4) the seed $z$ must remain masked throughout, since it is a long-term secret that persists across all decapsulations.

A fault injection that forces the comparison to always return ``equal'' converts ML-KEM into a CCA-vulnerable scheme: the attacker can submit malformed ciphertexts and always receive the real shared secret, enabling an adaptive chosen-ciphertext attack.  The countermeasure is temporal redundancy on the re-encryption and comparison, with independent comparison logic that would require two simultaneous faults to bypass.

With ML-KEM's compact modulus and linear NTT, first-order masking is both effective and affordable.  ML-DSA, as we shall see, is a substantially harder problem.


\section{ML-DSA: Attack Surfaces and Hardening}
\label{sec:mldsa-hardening}

ML-DSA\index{ML-DSA!hardware security} (FIPS 204~\cite{FIPS204}) operates over $\Zq$ with $q = 8380417$ (23 bits)---nearly double the bit-width of ML-KEM.  But the larger modulus is not the primary challenge.  ML-DSA introduces three complications absent from ML-KEM: a \emph{structural timing side channel} from rejection sampling, \emph{non-linear rounding operations} that resist masking, and a \emph{deterministic mode} that gives the attacker unlimited identical traces.

\subsection{The Rejection Sampling Timing Channel}
\label{subsec:rejection-sampling}

ML-DSA signing uses rejection sampling\index{rejection sampling!timing channel}: the signer generates a random masking vector $\mathbf{y}$, computes the signature candidate $\mathbf{z} = \mathbf{y} + c\mathbf{s}_1$ (where $c$ is the challenge and $\mathbf{s}_1$ is the secret key), and checks whether $\|\mathbf{z}\|_\infty < \gamma_1 - \beta$.  If the check fails, the signer restarts with fresh $\mathbf{y}$.  The number of iterations required depends on the secret $\mathbf{s}_1$: keys with larger coefficients cause more rejections.

For ML-DSA-65, the expected number of iterations is 4--7, with significant variance.  The attacker does not need to determine the exact count.  The total signing time is proportional to the iteration count and is trivially measurable---even remotely, via network timing.  With $10^4$--$10^5$ timed signatures, statistical analysis can distinguish keys.

The vulnerability is worth seeing in pseudocode, because it looks deceptively innocuous:

\begin{algorithm}[h]
\caption{ML-DSA signing with timing leak\index{ML-DSA!timing leak}}
\label{alg:mldsa-timing-leak}
\begin{algorithmic}[1]
\Require Secret key $(\mathbf{s}_1, \mathbf{s}_2)$, message $\mu$
\State $\text{attempts} \gets 0$
\Repeat
    \State $\text{attempts} \gets \text{attempts} + 1$
    \State Sample masking vector $\mathbf{y} \gets D_{\gamma_1}^{\ell}$
    \State Compute $\mathbf{w} \gets \mathbf{A}\mathbf{y}$
    \State Derive challenge $c \gets H(\mu \| \mathbf{w}_1)$
    \State Compute $\mathbf{z} \gets \mathbf{y} + c\,\mathbf{s}_1$ \Comment{Secret enters here}
    \State Check $\|\mathbf{z}\|_\infty < \gamma_1 - \beta$ \Comment{Rejection depends on $\mathbf{s}_1$}
\Until{all conditions satisfied}
\State \Return $(\mathbf{z}, \mathbf{h}, c)$ \Comment{Total time $\propto$ \texttt{attempts}}
\end{algorithmic}
\end{algorithm}

The variable \texttt{attempts}\index{rejection sampling!iteration count} is the timing channel.  Keys with larger coefficients in $\mathbf{s}_1$ cause more rejections, increasing the iteration count.  The attacker's procedure is straightforward: collect timing measurements over $10^4$--$10^5$ signatures, compute per-message timing statistics, correlate timing variation with known message structure, and recover secret-key bits through the resulting bias.  The attack is practical---it requires patience, not sophistication.

This is not a bug in the implementation; it is a structural property of the algorithm.  The only countermeasure is \emph{constant-time padding}\index{constant-time padding}: always execute $N_{\max}$ iterations, discarding the extra results.  The value $N_{\max}$ must be calibrated so that $\Pr[\text{iterations} > N_{\max}] < 2^{-128}$.  All iterations---including the discarded ones---must write to the output buffer unconditionally.  Only the valid result is selected at the end via a masked multiplexer.

The constant-time padding has a severe performance cost.  For ML-DSA-65 with $p \approx 0.2$ success probability per iteration, the geometric distribution gives $\Pr[\text{iterations} > N] = (1-p)^N = 0.8^N$.  To achieve $0.8^{N_{\max}} < 2^{-128}$, we need $N_{\max} > 128 / \log_2(1/0.8) \approx 398$ iterations.  Since the expected number of iterations is only $1/p = 5$, the worst-case latency exceeds the average by a factor of roughly 80.  This is the price of constant-time execution for a rejection-sampling-based scheme, and it is one of the principal engineering arguments for preferring ML-KEM (which has no rejection sampling) when latency-sensitive hardware is a priority.

\subsection{Deterministic vs.\ Hedged Mode}
\label{subsec:det-vs-hedged}

In deterministic mode\index{deterministic mode}, the randomness parameter \texttt{rnd} is set to zero.  The same message and secret key produce identical intermediate values across executions.  For DPA, this is the worst-case scenario: the attacker can average unlimited traces of the same computation, driving the noise to zero.

In \emph{hedged mode}\index{hedged mode}, \texttt{rnd} is drawn fresh from the DRBG for each signature.  Different executions with the same message now produce different intermediate values, preventing trace averaging.  For hardware implementations with access to a high-quality random number generator, hedged mode is strictly superior from a side-channel perspective.

\begin{svgraybox}
\textbf{Recommendation.}  Hardware implementations with a DRBG backed by a physical entropy source should implement hedged mode only.  Omitting deterministic mode reduces the certification analysis scope significantly and eliminates the DPA worst case.  If deterministic mode must be offered (e.g., for interoperability testing), it requires higher-order masking or a formal bound on the number of signatures permitted with the same message.
\end{svgraybox}


\subsection{Non-Linear Rounding and Norm Checks}
\label{subsec:rounding-norm}

ML-DSA's \texttt{HighBits} and \texttt{LowBits}\index{HighBits}\index{LowBits} functions decompose a polynomial coefficient into high-order and low-order parts via integer division and rounding.  These operations are non-linear and difficult to mask, especially for the non-power-of-2 modulus $q = 8380417$.  The transition from the masked to unmasked domain immediately before rounding is a concentrated leakage point.  The practical approach is \emph{late unmasking}: maintain the arithmetic masking as deep into the computation as possible, converting to Boolean representation only at the final moment before the non-linear rounding, and minimising the temporal window during which the unmasked value exists.

The infinity-norm check ($\|\mathbf{z}\|_\infty < \gamma_1 - \beta$) must evaluate all 256 coefficients without early exit.  An early exit---stopping as soon as one coefficient exceeds the threshold---reveals which coefficient failed, leaking information about the secret.  In the masked domain, each norm comparison requires A2B conversion (23 bits per coefficient), making this one of the most expensive masked operations in ML-DSA.


\subsection{Fault Injection on SampleInBall}
\label{subsec:fi-sampleinball}

The function \texttt{SampleInBall} generates the challenge polynomial $c$ with exactly $\tau$ non-zero coefficients (each $\pm 1$).  A fault that forces $c = 0$ reduces the signature equation to $\mathbf{z} = \mathbf{y}$: the output is the nonce itself.  Two such faults with different (real) challenges allow the attacker to compute $c_1 \mathbf{s}_1 - c_2 \mathbf{s}_1 = \mathbf{z}_1 - \mathbf{z}_2$, recovering the secret key algebraically.

The countermeasure is an integrity check on $c$: after sampling, verify that exactly $\tau$ coefficients are non-zero and that their values are in $\{-1, +1\}$.  This check is cheap (a single pass over 256 coefficients) and catches the most dangerous fault scenario.  Redundant encoding of the rejection loop counter adds a second layer of defence.


\subsection{The 23-Bit Cost Scaling Problem}
\label{subsec:23bit-cost}

Every masked operation in ML-DSA inherits a cost penalty from the 23-bit modulus.  Compared to ML-KEM's 12-bit $q = 3329$:
\begin{itemize}
\item Multiplier area: $\sim$$4\times$ (area scales roughly as the square of bit-width for schoolbook multiplication).
\item A2B/B2A conversion: $\sim$$2\times$ (cost is linear in bit-width: 23 vs.\ 12 masked carry-chain steps).
\item BRAM for masked polynomials: $\sim$$2\times$ (23-bit shares vs.\ 12-bit shares).
\end{itemize}

This cost scaling limits hardware reuse between ML-KEM and ML-DSA.  A unified PQC accelerator must either provision the wider datapath (wasting area when running ML-KEM) or maintain separate arithmetic units (wasting area when only one algorithm is active).

\subsection{Hint Generation and Secret Leakage}
\label{subsec:hint-generation}

The hint vector $\mathbf{h}$ in ML-DSA is a public output---it is included in the signature.  However, its \emph{generation} involves operations on the secret vectors $\mathbf{s}_2$ and $\mathbf{t}_0$.  The function \texttt{MakeHint} computes whether adding a correction changes the high-order bits of a value, which requires evaluating a rounding condition on secret-derived data.  The hint itself does not leak the secret (it is designed not to), but the \emph{process} of computing it---the power consumption, the timing of conditional evaluations---may leak through a side channel.

The practical impact is that hint computation must be treated as a secret operation even though its output is public.  The masked domain must extend through the hint generation, with unmasking deferred until the hint is written to the output buffer.

\subsection{Countermeasure Summary for ML-DSA}
\label{subsec:mldsa-countermeasures}

The complete hardening of ML-DSA requires:
\begin{itemize}
\item Constant-time rejection loop with $N_{\max}$ padding and unconditional buffer writes.
\item Hedged mode as the default (or only) signing mode.
\item DOM-1 Keccak for all SHAKE invocations on secret data.
\item Arithmetic masking (23-bit) on all operations involving $\mathbf{s}_1$, $\mathbf{s}_2$, and nonce $\mathbf{y}$.
\item Late unmasking for \texttt{HighBits}/\texttt{LowBits}, with minimal temporal exposure.
\item Full-vector norm evaluation without early exit.
\item Integrity verification of the challenge polynomial $c$ after \texttt{SampleInBall}.
\item Redundant encoding of the rejection loop counter.
\item Shuffling of NTT coefficient processing order.
\end{itemize}

The combined overhead is substantially higher than for ML-KEM: roughly 40--60\% additional latency and 3--4$\times$ additional area for the masked arithmetic, dominated by the 23-bit multipliers and the constant-time rejection loop padding.


\section{SLH-DSA: Attack Surfaces and Hardening}
\label{sec:slhdsa-hardening}

SLH-DSA\index{SLH-DSA!hardware security} (FIPS 205~\cite{FIPS205}) is architecturally the opposite of ML-KEM and ML-DSA.  There is no lattice algebra, no NTT, no polynomial multiplication.  All cryptographic operations reduce to invocations of Keccak/SHAKE.  This makes the security surface \emph{homogeneous}---the Keccak core is the only target---but the attack volume is enormous.

\subsection{The Sheer Volume Problem}
\label{subsec:slhdsa-volume}

A single SLH-DSA-128f signature requires approximately 35,000 Keccak invocations.  Many of these process secret data: the WOTS+ chain computations iterate $\text{value}_{i+1} = H(\text{value}_i)$, where each intermediate value is a secret derived from the master seed \texttt{sk\_seed}.  The key derivation step $\text{PRF}(\texttt{sk\_seed}, \text{address}_i)$ is the ideal CPA scenario: a fixed secret (\texttt{sk\_seed}), a known varying input (the public address), and many traces.

For a key used to sign 1000 messages, the attacker accumulates roughly $35 \times 10^6$ Keccak traces over related secrets.  If the per-trace SNR is moderate (as on an FPGA), first-order masking may be insufficient: second-order statistical attacks over millions of traces become feasible.

\subsection{WOTS+ Chains as CPA Targets}
\label{subsec:wots-cpa}

Each WOTS+ chain\index{WOTS+!chain computation} computes a sequence of hash values: starting from a secret leaf value derived from \texttt{sk\_seed}, each step applies Keccak to produce the next value in the chain.  Successive values are functionally related---each is the hash of the previous one---giving the attacker a ``related-input'' advantage beyond standard CPA.  The chain values that are revealed in the signature provide known plaintexts for verification of partial attacks.

This is qualitatively different from attacking ML-KEM's NTT, where each decapsulation processes an independent ciphertext.  In SLH-DSA, the secret material is structurally correlated across invocations, and the attacker benefits from this correlation.

\subsection{Masking Order: DOM-1 vs.\ DOM-2}
\label{subsec:masking-order}

The decision between first-order masking (DOM-1) and second-order masking (DOM-2) for SLH-DSA's Keccak should be driven by a quantitative evaluation:
\begin{equation}
\label{eq:masking-order-criterion}
\text{traces per key} \times \text{SNR}^{2d} \quad\text{vs.}\quad \text{attack complexity at order } d.
\end{equation}

If the product targets high-volume signing (e.g., a server signing thousands of TLS certificates per hour), the trace budget $N = 35{,}000 \times \text{signatures}$ may exceed the first-order attack threshold.  In such cases, DOM-2 for SLH-DSA's Keccak is justified even if DOM-1 suffices for ML-KEM and ML-DSA.  DOM-2 roughly cubes the Keccak area and requires $2 \times 1600 = 3{,}200$ fresh random bits per round (76,800 per invocation)---a significant but not prohibitive cost for a well-resourced hardware implementation.

\subsection{Fault Injection on Address and Leaf Index}
\label{subsec:slhdsa-fi}

Forcing the address register\index{address register (SLH-DSA)} to repeat a value causes the same WOTS+ leaf to be used for two different messages.  Each signature reveals a subset of the chain values; two signatures from the same leaf reveal a larger subset, potentially enabling forgery.  The countermeasure is redundant encoding of the address register (dual rail or ECC) with illegal-value detection and abort.

Complete signature buffering before emission is also required: SLH-DSA signatures are large (7--50~KB depending on the parameter set), and streaming partial signatures before the full computation completes would expose partial authentication paths without integrity guarantees.

\subsection{Countermeasure Summary}
\label{subsec:slhdsa-countermeasures}

SLH-DSA hardening is conceptually simple---protect Keccak, protect the address logic, buffer outputs---but the scale is formidable.  The masking randomness demand at DOM-1 is $35{,}000 \times 38{,}400 \approx 1.3 \times 10^9$ bits per signature, or roughly 170~MB.  At DOM-2, this doubles.  The entropy source must be dimensioned accordingly, and the randomness delivery path must itself be protected against fault injection.

An additional countermeasure specific to SLH-DSA is \emph{shuffling of independent hash chain computations}\index{shuffling!WOTS+ chains} within a WOTS+ signature.  A WOTS+ signature consists of $\ell$ independent chains (where $\ell$ depends on the parameter set, typically 35--67).  The order in which these chains are computed can be randomised without affecting the output.  This prevents the attacker from aligning traces to specific chains, multiplying the required number of traces by $\ell$.

Key generation deserves special attention: it computes the full hypertree, requiring millions of Keccak invocations over \texttt{sk\_seed}.  This is a massive SCA trace volume concentrated in a single operation.  The recommendation is to perform key generation in a controlled environment (e.g., during manufacturing, with physical shielding) or with an elevated masking order.

SLH-DSA is the first PQC algorithm certifiable with SCA resistance once a DOM-protected Keccak core is available, precisely because its security surface is homogeneous.  There are no lattice-specific complications, no domain conversions, no rejection sampling---just Keccak, at scale.


\section{Cross-Component Interactions}
\label{sec:cross-component}

The previous three sections analysed each PQC algorithm in isolation.  But a real hardware cryptographic module does not contain a single algorithm in a vacuum.  It contains an entropy source, a DRBG, multiple PQC cores, shared bus interconnects, and a finite-state machine that orchestrates them.  The interfaces between these components introduce systemic risks that only manifest at the integration level.  Testing components individually and declaring each ``secure'' provides no guarantee about the integrated system---the interactions themselves are the attack surface.

\subsection{The DRBG as Systemic Vulnerability}
\label{subsec:drbg-systemic}

The DRBG's internal state (1600 bits for a Keccak-based XDRBG~\cite{KampanakisDang2023}) is more valuable than any individual PQC key: compromising the state yields all keys, nonces, and masks generated until the next reseed.  The state is a long-lived secret processed repeatedly by Keccak, and each invocation provides the attacker with a correlated trace.  Between reseeds, if the DRBG services $K$ requests, the attacker has $K$ traces of the same (evolving) state.  Without masking, CPA recovers the state; with DOM-1, the attacker needs second-order analysis across the $K$ traces.

Differential Fault Analysis (DFA)\index{Differential Fault Analysis} on the Keccak permutation is equally dangerous: a single-byte fault in rounds 22--24 of Keccak-$f[1600]$ enables full state recovery with as few as 2--4 faults~\cite{BagheriEtAl2015}.  This is a more direct attack than CPA and requires fewer interactions, making fault-injection countermeasures on the DRBG's Keccak core at least as important as masking.

The DRBG's FSM also presents a target.  A fault that causes a reseed to be interpreted as a generate operation means the state is not refreshed---potentially yielding correlated outputs.  The reverse fault (generate interpreted as reseed) may overwrite state with unprocessed data.  Redundant FSM encoding with trap-state detection mitigates both scenarios.

\subsection{The Reseed Integrity Problem}
\label{subsec:reseed-integrity}

The entropy path from the physical random number generator to the DRBG crosses a physical bus.  If a fault injection corrupts the entropy bits in transit, the DRBG absorbs corrupted data as if it were genuine entropy.  This is \emph{worse} than no reseed: if the corruption is attacker-controlled, the attacker can steer the DRBG state to a known value, compromising all subsequent outputs.

The countermeasure is an integrity check on the reseed path: a CRC or MAC computed over the entropy block, verified by the DRBG before absorption, with abort and zeroisation on mismatch.  The integrity check itself must be protected against fault injection---a redundant computation with independent comparison logic.

\subsection{The Circular Dependency}
\label{subsec:circular-dependency}

Consider the following dependency chain\index{circular dependency (masking randomness)}:
\begin{enumerate}
\item The DRBG (e.g., XDRBG based on SHAKE-256) produces randomness for PQC algorithms.
\item The DRBG's Keccak core must be masked (DOM-1 or DOM-2) to prevent state recovery via CPA.
\item Masking Keccak requires fresh random bits---38,400 per invocation at DOM-1.
\item Those fresh random bits must come from... the DRBG?
\end{enumerate}

This is circular: the DRBG cannot produce the randomness needed to protect itself.  The solution is to break the cycle by sourcing the masking randomness for the DRBG's Keccak \emph{directly from the physical entropy source}---raw bits that have passed the health test but have not yet been conditioned by the DRBG.  These bits are not cryptographically conditioned, but they do not need to be: masking randomness requires statistical uniformity, not computational unpredictability.  The entropy source's raw output, post-health-test, is sufficient.

This creates a clear hierarchy: the entropy source feeds raw randomness to the DRBG's masking logic (bypassing the DRBG), while the DRBG feeds conditioned randomness to the PQC algorithms' masking logic and operational randomness needs.

The architecture has three layers, each with a different randomness quality:
\begin{enumerate}
\item \textbf{Raw entropy} (post-health-test, pre-conditioning): feeds DRBG masking logic.  Needs statistical uniformity but not cryptographic unpredictability.
\item \textbf{Conditioned entropy} (DRBG output): feeds PQC algorithm masking logic and shuffling.  Needs computational unpredictability.
\item \textbf{Algorithm-level randomness} (also DRBG output, domain-separated): feeds PQC key generation, nonces, encapsulation coins.  Needs full cryptographic quality with per-consumer domain separation.
\end{enumerate}

The domain separation at layer 3 is essential for composability: the security proof of each PQC algorithm assumes independent randomness.  If two algorithms share a DRBG without domain separation (i.e., without distinct personalisation strings), a correlation between their random inputs may exist, and the composition is not provably secure~\cite{FIPS203, FIPS204}.

\subsection{Shared Keccak Core Contention}
\label{subsec:shared-keccak}

In area-constrained designs, a single Keccak core may be shared among the DRBG, ML-KEM, ML-DSA, and SLH-DSA.  Sharing creates two problems.  First, \emph{contention}\index{Keccak!shared core contention}: when two modules need Keccak simultaneously, one must wait, and the wait time depends on the other module's secret-dependent computation---a timing side channel.  Second, \emph{state residue}: if the Keccak core's registers are not zeroised between invocations from different security domains, secret state from one domain may leak into the power trace of another domain's computation.

The countermeasures are either separate Keccak instances per security domain (expensive in area) or a constant-time arbiter\index{constant-time arbiter} that services requests in a fixed schedule regardless of demand, combined with mandatory zeroisation between context switches.

The choice between these approaches is an engineering trade-off with certification implications.  Separate instances provide stronger isolation (each domain has no way to influence another's timing) but require more silicon area and more masking randomness.  A shared instance with a constant-time arbiter is more area-efficient but requires a careful security argument documenting why the arbiter's scheduling policy does not create cross-domain timing channels.  The security argument must address not just the Keccak invocations themselves but also the queuing delays when multiple requests arrive simultaneously.

\subsection{Frequent Reseeding as a Side-Channel Countermeasure}
\label{subsec:reseed-sca}

An underappreciated advantage of high-throughput entropy sources is their impact on side-channel resistance.  If the DRBG reseeds every $K$ invocations, the attacker has at most $K$ traces per DRBG state.  With DOM-1 Keccak requiring approximately $10^6$ traces for a successful second-order attack, setting $K \ll 10^6$ ensures the state changes before the attacker accumulates enough traces.

This is a quantifiable, multiplicative security gain.  A DRBG reseeded every 1000 invocations from a high-throughput entropy source has an effective SCA resistance that compounds with the masking order: the attacker needs traces \emph{per state}, but each state lasts only $K$ invocations.  This turns the reseed frequency into a formal security parameter, with a documented bound: ``maximum $K$ invocations between reseeds, where $K$ is chosen such that $K \ll$ traces needed for order-$d$ attack.''

\begin{example}[Reseed Frequency Dimensioning]
\label{ex:reseed-frequency}
Suppose a DOM-1 masked Keccak requires $10^6$ traces for a successful second-order CPA attack (a reasonable estimate for moderate SNR on an FPGA).  If the DRBG reseeds every $K = 1000$ invocations, the attacker has 1000 traces per state---three orders of magnitude below the attack threshold.  Even if the attacker collects traces across $10^3$ different states, each batch is independent and cannot be combined.  The reseed frequency effectively reduces the attacker's useful trace budget by a factor of $1000\times$.  A physical entropy source capable of delivering a new 256-bit seed every 10~ms (25.6~Kbit/s---trivial for any hardware RNG) supports this reseed frequency at the cost of one additional Keccak invocation per reseed.
\end{example}

\subsection{Comparative Hardening Cost}
\label{subsec:comparative-cost}

Table~\ref{tab:comparative-hardening} summarises the hardening cost across the three algorithms.  The differences are stark: ML-KEM is the cheapest to protect, SLH-DSA is the simplest (only Keccak), and ML-DSA is the most complex (wider datapath, rejection sampling, rounding).

\begin{table}[t]
\caption{Comparative hardening cost for first-order masking across NIST PQC standards}
\label{tab:comparative-hardening}
\begin{tabular}{p{3cm}p{3cm}p{3cm}p{3cm}}
\hline\noalign{\smallskip}
\textbf{Aspect} & \textbf{ML-KEM} & \textbf{ML-DSA} & \textbf{SLH-DSA} \\
\noalign{\smallskip}\svhline\noalign{\smallskip}
Modulus bit-width & 12 bits & 23 bits & N/A (hash-only) \\
NTT masking cost & Free (linear) & Free (linear) & N/A \\
A2B/B2A cost & ${\sim}12$ cycles & ${\sim}23$ cycles & N/A \\
Keccak invocations & ${\sim}6$ per decaps & ${\sim}10$ per sign & ${\sim}35{,}000$ per sign \\
DOM-1 randomness & ${\sim}230$ Kbit & ${\sim}384$ Kbit & ${\sim}1.3$ Gbit \\
Timing channel & None (structural) & Rejection sampling & None (structural) \\
Latency overhead & ${\sim}20$--$30\%$ & ${\sim}40$--$60\%$ & ${\sim}20$--$40\%$ \\
Area overhead & ${\sim}2\times$ multiplier & ${\sim}4\times$ multiplier & ${\sim}2\times$ Keccak \\
\noalign{\smallskip}\hline
\end{tabular}
\end{table}

The cross-component interactions analysed in this section---reseed integrity, circular dependencies, shared resource contention, and reseed frequency---must all be documented in the security target for any hardware certification.  We now turn to the certification frameworks themselves.


\section{The Certification Landscape}
\label{sec:certification}

A hardened implementation is only as valuable as the certification it can achieve.  Regulated industries---government, defence, finance, critical infrastructure---require cryptographic modules to be independently evaluated and certified against recognised standards.  This section surveys the three major certification frameworks and maps the security requirements from the previous sections to specific framework clauses.

\subsection{FIPS 140-3}
\label{subsec:fips140}

The \emph{Federal Information Processing Standard 140-3}\index{FIPS 140-3}~\cite{FIPS140-3} defines security requirements for cryptographic modules at four levels of increasing stringency.  It is mandatory for all cryptographic modules used by US federal agencies and widely adopted by private industry as a baseline.

\begin{definition}[FIPS 140-3 Security Levels]
\label{def:fips-levels}
\begin{itemize}
\item \textbf{Level 1:} Functional correctness.  Algorithm validation via Known Answer Tests (KATs) and the Automated Cryptographic Validation Protocol (ACVP).  No physical security requirements.
\item \textbf{Level 2:} Role-based authentication plus evidence of tamper (tamper-evident coatings or seals).  Minimal physical security.
\item \textbf{Level 3:} Tamper resistance (active response to physical penetration) and \emph{non-invasive side-channel testing} per ISO 17825~\cite{ISO17825}.  This is the level at which SCA countermeasures become mandatory.
\item \textbf{Level 4:} Complete envelope of physical protection.  Environmental failure protection and formal verification methods.  Few commercial products target Level~4.
\end{itemize}
\end{definition}

For PQC hardware, the critical transition is from Level 2 to Level 3: the moment SCA testing becomes mandatory.  A Level 3 evaluation requires the module to resist non-invasive attacks (power analysis, EM analysis) as tested by an accredited laboratory.  The masking, shuffling, and fault countermeasures developed in this chapter are precisely what Level 3 demands.

\subsection{Common Criteria}
\label{subsec:common-criteria}

The \emph{Common Criteria}\index{Common Criteria} (CC)~\cite{CommonCriteria2022} is an international framework (ISO/IEC 15408) that evaluates products against a \emph{Protection Profile} (PP) specifying the required security functionality.  Evaluation Assurance Levels (EAL) range from EAL1 (functionally tested) to EAL7 (formally verified design and tested).  For cryptographic hardware, the critical metric is the \emph{vulnerability assessment} level:

\begin{itemize}
\item \textbf{AVA\_VAN.3:}  Resistance to attackers with ``enhanced-basic'' attack potential.  Basic SCA testing.
\item \textbf{AVA\_VAN.4:}  Resistance to attackers with ``moderate'' attack potential.  Laboratory-grade SCA testing with power and EM analysis is mandatory.
\item \textbf{AVA\_VAN.5:}  Resistance to attackers with ``high'' attack potential.  Includes fault injection testing.
\end{itemize}

Products targeting the European market (particularly government and payment systems) typically require EAL4+ with AVA\_VAN.4 or AVA\_VAN.5.  The CC evaluation is more flexible than FIPS 140-3 in that the Protection Profile can be tailored, but also more demanding at the highest levels: AVA\_VAN.5 explicitly requires resistance to fault injection attacks, which FIPS 140-3 Level 3 addresses only indirectly through physical tamper resistance.

\subsection{BSI Guidelines}
\label{subsec:bsi}

The German Federal Office for Information Security (BSI)\index{BSI} publishes Technical Reports that often set the de facto standard for European evaluations.  For hardware security:

\begin{itemize}
\item \textbf{TR-03158:}  Side-channel analysis resistance evaluation methodology~\cite{BSI-TR-03158}.  Specifies the measurement setup, statistical tests, and pass/fail criteria.
\item \textbf{AIS 20/31:}  Evaluation of random number generators.  Defines Physical True Random Number Generator classes (PTG.1, PTG.2, PTG.3) with increasing requirements on the stochastic model, health testing, and conditioning.
\end{itemize}

BSI requirements are roughly aligned with CC but often more prescriptive in methodology.  For a PQC hardware module targeting both FIPS 140-3 Level 3 and Common Criteria EAL4+/AVA\_VAN.4, the BSI TR-03158 methodology provides the most detailed blueprint for the SCA evaluation.

\subsection{Certification Mapping}
\label{subsec:cert-mapping}

\begin{table}[t]
\caption{Mapping security requirements to certification framework clauses}
\label{tab:cert-mapping}
\begin{tabular}{p{2.5cm}p{3.5cm}p{3.5cm}p{2.5cm}}
\hline\noalign{\smallskip}
\textbf{Requirement} & \textbf{FIPS 140-3} & \textbf{Common Criteria} & \textbf{BSI} \\
\noalign{\smallskip}\svhline\noalign{\smallskip}
Algorithm correctness & KATs + ACVP (SP 800-140D) & ADV\_FSP + ATE\_COV & AIS 20/31 (RNG) \\
SCA resistance & ISO 17825 (Level 3+) & AVA\_VAN.4/.5 & TR-03158 \\
FI resistance & Physical security (Level 3+) & AVA\_VAN.5 & TR-03158 \\
Zeroisation & SP 800-140D & FCS\_CKM.4 & CC-aligned \\
RNG quality & SP 800-90A/B/C & FCS\_RBG & AIS 20/31 \\
Formal verification & Expected at Level 4 & ADV\_SPM (EAL5+) & EAL5+ \\
\noalign{\smallskip}\hline
\end{tabular}
\end{table}


\subsection{The Practical Path}
\label{subsec:practical-path}

For a vendor developing a PQC hardware module, the certification process involves several stages.  First, algorithm validation: each PQC algorithm must pass NIST's ACVP testing (for FIPS) or equivalent functional testing (for CC).  Second, documentation: the security target (CC) or security policy (FIPS) must describe all cryptographic boundaries, key management, and security mechanisms.  Third, laboratory evaluation: an accredited lab performs SCA testing (power, EM, sometimes laser fault injection) against the claimed security level.  Fourth, report and certification body review.

\begin{svgraybox}
\textbf{Timeline and cost.}  A FIPS 140-3 Level 3 certification typically takes 12--24 months from submission to certificate, depending on laboratory queue length and the number of findings that require remediation.  Costs range from \$200K to \$500K for laboratory fees, plus internal engineering effort.  Common Criteria EAL4+/AVA\_VAN.4 is comparable in duration but often more expensive due to the depth of documentation required (the Security Target alone can exceed 200 pages).  Delta certifications---updating an existing certificate for a new algorithm (e.g., adding ML-KEM to an already-certified module)---are faster and cheaper, typically 6--12 months.
\end{svgraybox}

The PQC-specific complication is that certification laboratories are still developing their testing methodologies for lattice-based and hash-based algorithms.  The attack libraries used for SCA testing of AES and RSA do not directly apply to NTT-based schemes.  Early movers in PQC hardware certification face the challenge---and the advantage---of helping to define the evaluation methodology.

\subsection{What Certification Does Not Cover}
\label{subsec:cert-gaps}

It is worth being explicit about the limits of certification.  FIPS 140-3 and Common Criteria certify a \emph{specific version} of a \emph{specific product} against a \emph{specific set of requirements}.  They do not guarantee security in perpetuity.  New attack techniques---discovered after certification---may invalidate the security claims.  The certification is a snapshot, not a proof.

For PQC specifically, several open questions remain in the certification landscape as of 2026:
\begin{itemize}
\item \textbf{NTT-specific SCA test vectors:} Certification labs have decades of experience testing AES and RSA/ECC implementations.  Standardised test procedures for NTT-based leakage detection are still being developed.
\item \textbf{Rejection sampling evaluation:} How to evaluate the timing channel from ML-DSA's variable-time signing in a systematic, reproducible way.
\item \textbf{Hash-based signature volume:} Whether the SCA evaluation for SLH-DSA should account for the cumulative trace volume over the key's expected lifetime, not just a single operation.
\item \textbf{DRBG-PQC composition:} Whether the DRBG and PQC modules can be certified separately or must be evaluated as a composed system.
\end{itemize}

Beyond implementation attacks, there is a subtler gap: certification fixes a \emph{security level} at the time of evaluation, but the underlying hardness assumptions erode as lattice-reduction algorithms\index{lattice reduction!progressive improvement} improve.  Table~\ref{tab:lattice-reduction-evolution} traces the progression of best-known lattice sieving complexities over a fifteen-year span.  The exponent has dropped from $2^{0.292d}$ to $2^{0.280d}$---a modest-looking shift that translates to a 5--12~bit loss in concrete security for typical lattice dimensions.

\begin{table}[t]
\caption{Evolution of lattice reduction algorithms\index{lattice reduction!algorithm evolution}: best-known asymptotic sieving complexity as a function of lattice dimension~$d$}
\label{tab:lattice-reduction-evolution}
\begin{tabular}{p{1.5cm}p{3.5cm}p{2.5cm}p{4.5cm}}
\hline\noalign{\smallskip}
\textbf{Year} & \textbf{Algorithm} & \textbf{Complexity} & \textbf{Impact} \\
\noalign{\smallskip}\svhline\noalign{\smallskip}
2010 & BKZ 2.0 & $2^{0.292d}$ & Established baseline for NIST parameter selection \\
2016 & Progressive BKZ & $2^{0.290d}$ & Minor improvement; parameters still comfortable \\
2018 & G6K Sieving & $2^{0.285d}$ & 5--10 bit security reduction at $d \approx 800$ \\
2021 & BDGL Sieving & $2^{0.283d}$ & Continued erosion; no parameter changes needed \\
2024 & Tensor methods & $2^{0.280d}$ & Under active investigation \\
\noalign{\smallskip}\hline
\end{tabular}
\end{table}

The NIST ML-KEM and ML-DSA parameters were chosen with a substantial margin above the $2^{0.292d}$ baseline, so the current erosion does not demand parameter changes.  But the trend is clear and monotonic: every few years, the exponent drops.  A certified module deployed in 2026 with a 20-year operational lifetime will face lattice-reduction algorithms that do not yet exist.  Certification captures the security posture at a point in time; periodic reassessment against updated hardness estimates is essential for any long-lived deployment.

These gaps mean that a vendor pursuing early PQC hardware certification must engage proactively with the evaluation laboratory, proposing test methodologies rather than waiting for published standards.  The technical analysis in this chapter provides the foundation for those proposals.

The interplay between hardware security and certification creates a virtuous cycle: the threat analysis drives the countermeasure design, the countermeasure design structures the certification evidence, and the certification process validates (or challenges) the threat analysis.  A well-designed PQC hardware module treats security analysis and certification planning as parallel activities, not sequential ones.

The next chapter addresses a different but equally practical challenge: how organisations migrate their existing cryptographic infrastructure to post-quantum algorithms.  Where this chapter asked ``how do we build a secure PQC implementation?'', the next asks ``how do we deploy it into a world that still runs on RSA and elliptic curves?''  The hardware security analysis developed here---particularly the certification requirements and the cross-component interaction model---will reappear when we discuss the constraints that certified hardware modules impose on migration timelines and hybrid protocol design.


\section*{Problems}
\addcontentsline{toc}{section}{Problems}

\begin{prob}
\label{prob:ch13-cpa}
\textbf{CPA walkthrough.}\\
An attacker targets a single coefficient $b$ of the ML-KEM secret polynomial $\hat{s}$, where $q = 3329$.  The targeted NTT butterfly computes $a' = a + \omega b \bmod q$ with public twiddle factor $\omega = 17$ and known input $a$.
\begin{enumerate}
\item[(a)] For each key hypothesis $k \in \{0, 1, \ldots, 3328\}$, write the expression for the predicted intermediate value $v(k) = a + 17k \bmod 3329$.
\item[(b)] Using the Hamming weight model, compute $\text{HW}(v(k))$ for $k = 0, 1, 2, 3$ and $a = 1000$.  Show that different hypotheses produce different predicted power values.
\item[(c)] Suppose the attacker collects $N = 5000$ traces with different values of $a$ (drawn from decapsulations with different ciphertexts).  For the correct hypothesis $k = b$, the predicted power values are correlated with the measured power at the targeted clock cycle.  Explain why incorrect hypotheses $k \neq b$ produce near-zero correlation as $N$ grows.
\item[(d)] If the SNR is 0.1 (typical for an FPGA), estimate the minimum number of traces needed for a successful first-order CPA using the rule of thumb $N \approx \alpha / \text{SNR}^2$ with $\alpha \approx 5$.  How does this change if the implementation uses first-order arithmetic masking?
\end{enumerate}
\end{prob}

\begin{prob}
\label{prob:ch13-masking-cost}
\textbf{Masking cost estimation.}\\
A hardware ML-KEM-768 decapsulation on FPGA runs at 100~MHz and takes 50~$\mu$s unmasked.
\begin{enumerate}
\item[(a)] The NTT involves $256 \times \log_2 256 / 2 = 1024$ butterfly operations.  Under first-order arithmetic masking, the NTT must be computed on two shares.  If the multiplier is not duplicated (i.e., shares are processed sequentially), what is the NTT's new latency?  If the multiplier is duplicated, what is the area cost?
\item[(b)] The decapsulation requires 6 Keccak invocations on secret data.  DOM-1 adds one pipeline cycle per Keccak round.  At 100~MHz, compute the additional latency for the 6 invocations.
\item[(c)] There are approximately 1536 A2B conversions (256 coefficients $\times$ 6 polynomials) during decapsulation.  At 12 cycles per conversion, compute the total conversion overhead.
\item[(d)] Sum the overheads from (a)--(c) and express the masked decapsulation latency as a percentage increase over the unmasked baseline.  Identify the dominant cost contributor.
\end{enumerate}
\end{prob}

\begin{prob}
\label{prob:ch13-slhdsa-randomness}
\textbf{SLH-DSA masking randomness budget.}\\
SLH-DSA-128f requires approximately 35,000 Keccak invocations per signature.
\begin{enumerate}
\item[(a)] At DOM-1, each Keccak invocation consumes 38,400 bits of fresh randomness.  Compute the total randomness per signature in bits and in megabytes.
\item[(b)] If the device signs at 10 signatures per second, what is the required randomness throughput in Mbit/s?
\item[(c)] A typical FPGA-based QRNG produces 100~Mbit/s of raw (post-health-test) output.  Is this sufficient for DOM-1?  For DOM-2 (which doubles the per-invocation randomness)?
\item[(d)] If the randomness source cannot keep up, the signing operation must stall.  Explain why this stall is itself a potential side channel and propose a mitigation.
\end{enumerate}
\end{prob}

\begin{prob}
\label{prob:ch13-rejection-timing}
\textbf{ML-DSA rejection sampling timing channel.}\\
Consider ML-DSA-65, where each signing attempt succeeds with probability $p \approx 0.2$.
\begin{enumerate}
\item[(a)] The number of attempts follows a geometric distribution.  Compute the expected number of attempts and the standard deviation.
\item[(b)] If each attempt takes 100~$\mu$s, and the attacker can measure total signing time with 1~$\mu$s precision, how many distinct timing bins can the attacker observe?
\item[(c)] Explain why constant-time padding to $N_{\max}$ iterations eliminates this channel.  Compute $N_{\max}$ such that $\Pr[\text{iterations} > N_{\max}] < 2^{-128}$.
\item[(d)] What is the worst-case signing latency with this padding, and how does it compare to the expected latency without padding?
\end{enumerate}
\end{prob}

\begin{prob}
\label{prob:ch13-certification}
\textbf{Certification mapping exercise.}\\
A vendor is developing a PQC hardware module that implements ML-KEM-768 and ML-DSA-65.  The module targets FIPS 140-3 Level 3 certification and Common Criteria EAL4+/AVA\_VAN.4.
\begin{enumerate}
\item[(a)] List the specific countermeasures (from Sects.~\ref{sec:countermeasure-foundations}--\ref{sec:mldsa-hardening}) that are mandatory to satisfy the SCA resistance requirements at these certification levels.
\item[(b)] The vendor's DRBG uses SHAKE-256 (XDRBG construction).  Which certification framework clause governs the DRBG's algorithm correctness?  Which governs its SCA resistance?
\item[(c)] The module shares a single Keccak core between the DRBG and both PQC algorithms.  Identify the cross-component risks (Sect.~\ref{sec:cross-component}) that must be documented in the Security Target and tested by the evaluation laboratory.
\item[(d)] If the vendor later adds SLH-DSA-128f support, what additional security analysis is required?  Can this be handled as a delta certification?
\end{enumerate}
\end{prob}

\begin{prob}
\label{prob:ch13-circular}
\textbf{The circular dependency.}\\
A PQC hardware module uses a DRBG with DOM-1 masked Keccak.  The masking requires 38,400 bits of fresh randomness per Keccak invocation.
\begin{enumerate}
\item[(a)] Explain why the DRBG cannot use its own output as masking randomness for its own Keccak core.
\item[(b)] The module has access to a QRNG that produces 200~Mbit/s of raw output.  The DRBG processes 100 Keccak invocations per second.  Compute the randomness consumed by the DRBG's masking and the fraction of the QRNG's output this represents.
\item[(c)] The remaining QRNG output is conditioned by the DRBG for use as masking randomness by the PQC modules.  If ML-KEM decapsulation requires 6 masked Keccak invocations and runs at 1000 decapsulations per second, compute the total masking randomness demand for ML-KEM's Keccak alone.  Is the conditioned output from the DRBG sufficient?
\item[(d)] Propose an architecture that resolves the circular dependency while ensuring that both the DRBG and PQC modules receive adequate masking randomness.
\end{enumerate}
\end{prob}
